\section{Literature Review}
\label{sec:literature_review}

This section reviews three streams of literature that form the foundation of our work. The first concerns service network design (SND) formulations, tracing the evolution from cost-minimization models with exogenous demand to revenue-maximizing models that endogenize passenger choice. The second concerns passenger choice models and their integration into optimization frameworks. The third concerns dynamic discretization and iterative refinement methods for time-space network optimization.

\subsection{Service Network Design}

Service network design encompasses a broad class of optimization problems in which a planner determines which services to operate, when to schedule them, and how to allocate resources across a network to satisfy transportation demand. Foundational work on transportation network design and freight service network design established this perspective by formulating service selection, routing, and resource allocation decisions as fixed-charge and multi-commodity network design problems, often represented on time-expanded or time-space networks \citep{magnanti1984network, crainic1986multicommodity, crainic2000service, andersen2009service}. These formulations capture the temporal dimension of service operations through discretized time-space graphs, where nodes represent locations at specific time points and arcs represent service movements or idle resource holding. The resulting mixed-integer programs must balance operating costs against service quality, subject to resource conservation, capacity, and scheduling constraints.

A central challenge in SND is the representation of demand. Classical models treat demand as exogenous---fixed independently of the service plan---and focus on cost minimization for a given set of origin-destination flows \citep{crainic2000service, andersen2009service, boland2017continuous}. In this stream, the demand side is usually represented as a set of fixed commodities, while the optimization model determines the service frequencies, schedules, capacities, and asset movements required to move those commodities at minimum cost. However, in passenger transportation settings, demand is inherently endogenous: travelers choose among available itineraries based on service attributes such as travel time, number of connections, and schedule convenience. This observation has motivated a substantial body of work on integrating discrete choice models into network optimization.

The airline industry has been a primary application domain for SND with endogenous demand. Early contributions by \cite{abara1989applying} and \cite{hane1995fleet} modeled the fleet assignment problem as large-scale integer programs, assuming exogenous demand and fixed revenues. Recognizing that leg-based demand models fail to capture the network effects of passenger routing, \cite{barnhart2002itinerary} shifted the modeling paradigm to itinerary-based demand, enabling a more accurate representation of how passengers flow through hub-and-spoke networks. Subsequent refinements incorporated enhanced revenue modeling \citep{barnhart2009airline} and frequency competition to account for congestion effects \citep{vaze2012modeling}.

Building on these modeling advances, \cite{wei2020airline} formulated the joint timetabling and fleet assignment problem as a discrete-time service design problem with embedded passenger choice, solving it via an iterative heuristic. While this work demonstrated the revenue benefits of endogenizing demand, applying the model to full-scale instances involving major carriers presents substantial computational challenges due to the scale of the resulting time-space networks. Addressing these tractability concerns, \cite{yan2022choice} introduced a decomposition-based solution approach that reduces dimensionality while preserving choice characteristics, demonstrating consistent improvements over benchmark methods in both computational efficiency and solution quality. Nevertheless, existing approaches remain heuristic or decomposition-based and do not provide an exact framework with provable global optimality for the large-scale, itinerary-level formulation considered in this paper.

Despite these advances, the SND literature with endogenous demand has focused predominantly on airline applications. Although related planning problems in public transit and bus networks have been extensively studied under transit network design, bus network design, and frequency-setting frameworks \citep{ceder1986bus, guihaire2008transit, farahani2013review}, these studies typically emphasize route design, frequency setting, scheduling, and passenger assignment rather than an exact integration of itinerary-level discrete choice with resource-constrained service network design. This gap motivates our focus on an exact approach for service network design with passenger choice that preserves the operational coupling between service scheduling, resource assignment, and endogenous demand.

\subsection{Passenger Choice Models}

A key challenge in service network design with endogenous demand is the tractable representation of passenger choice behavior within optimization models. Discrete choice models provide the theoretical foundation for predicting how travelers select among competing alternatives based on service attributes. The Multinomial Logit (MNL) model, derived from random utility maximization, is the most widely used specification due to its closed-form choice probabilities and its tractability in estimation and optimization \citep{mcfadden1974conditional, benakiva1985discrete}. Nested logit models extend the MNL framework by allowing alternatives within the same nest to share unobserved utility components, thereby relaxing the strict independence of irrelevant alternatives property for similar alternatives such as itineraries with common legs, connection patterns, or departure-time characteristics \citep{mcfadden1978modeling, benakiva1985discrete}.

A key breakthrough in tractably embedding choice behavior into optimization models came from choice-based linear programming formulations in network revenue management \citep{liu2008choice}, the Sales-Based Linear Program (SBLP) of \cite{gallego2015general}, and the attractiveness-based formulations of \cite{di2014attractiveness}. These frameworks provide linear representations of choice-consistent demand allocation, enabling their integration into mixed-integer optimization models without directly carrying the nonlinear choice probabilities of classical logit formulations. The SBLP linearization of the General Attraction Model (GAM) is particularly useful because it models substitution and demand recapture through attraction parameters while retaining a compact linear structure. In contrast to simple attraction models that may overestimate recapture when preferred alternatives are unavailable, the GAM introduces adjusted attractions and shadow attractions to moderate substitution effects \citep{gallego2015general}. This feature makes the SBLP/GAM framework well suited for service design problems in which the set of offered itineraries is itself an endogenous outcome of the schedule.

The integration of these choice models with service network design remains computationally challenging. Choice-based formulations require itinerary-level variables and constraints, and the number of feasible itineraries can grow combinatorially with the size of the time-space network, the number of allowed connections, and the temporal resolution of the planning horizon. As a result, even when the choice model admits a linear representation, the resulting mixed-integer program may remain extremely large and weakly relaxed. Our work builds on the SBLP/GAM modeling framework, but addresses this algorithmic bottleneck through itinerary aggregation and iterative refinement.

\subsection{Dynamic Discretization and Iterative Refinement}

A persistent challenge in both general SND and its passenger-transportation variants is the representation of time. Traditional models rely on time-space networks that require a discretization of the planning horizon into fixed intervals. Coarse discretization compromises model accuracy by conflating services that operate at different times within the same interval, while fine discretization leads to an explosion in problem size, rendering large-scale instances intractable.

To address these limitations, \cite{boland2017continuous} introduced Dynamic Discretization Discovery (DDD) for the Continuous-Time Service Network Design Problem (CTSNDP). The CTSNDP formulates service scheduling in a continuous time domain, avoiding the approximation errors inherent in static discretization. The DDD algorithm starts with a coarse discretization and iteratively refines the time-space network by adding only the time points required to verify optimality or improve the solution. This allows the solver to work with a significantly smaller network than would be required by a uniformly fine static discretization, while theoretically guaranteeing convergence to the optimal solution of the continuous-time problem. The DDD framework represents a paradigm shift from ``construct-then-solve'' to ``construct-while-solving,'' and has proven effective for freight service network design where demand is exogenous and the primary challenge is temporal compression of the network.

However, when passenger choice is endogenized, the computational bottleneck shifts from the physical network to the logical passenger dimension. The number of feasible itineraries grows combinatorially with the temporal resolution, and the choice-consistency constraints tightly couple demand allocation across markets through shared service arcs. DDD, by design, addresses only the physical network compression and does not provide a mechanism for managing this itinerary-level explosion. Decomposition approaches \citep{yan2022choice} partially address the scale issue but sacrifice the coupling between markets.

In this paper, we adopt the general principle of iterative discretization refinement but depart from the DDD framework in a fundamental way. Rather than treating discretization refinement as the sole driver of convergence, we introduce itinerary aggregation as a complementary dimension of relaxation. By constructing super-itineraries that compactly represent equivalence classes of original itineraries, our approach simultaneously compresses both the physical network and the passenger demand space. The resulting framework---which we term the Itinerary Aggregation and Dynamic Refinement (\algoNotation)---produces tight lower bounds, generates feasible upper bounds via a structure-aware repair algorithm, and guarantees finite convergence to the global optimum through targeted refinement of both time points and itinerary partitions.
